\documentclass[american]{cv}
\usepackage{pslatex}
\usepackage[T1]{fontenc}
\usepackage{geometry}
\geometry{verbose,letterpaper,tmargin=1in,bmargin=1in,lmargin=1in,rmargin=1in}
\usepackage{babel}

\makeatletter

\begin{document}

\leftheader{Georgia Tech Research Institute\\
8800 Redstone Gateway SW Suite 100\\
Huntsville, AL 35808}

\rightheader{tel: (256) 716-2158\\
jon.fox@drfox.com\\
jon.fox@gtri.gatech.edu}


\title{Dr. Jon R. Fox}

\maketitle
Applied Physicist with 25 years of post-doctoral experience spanning MEMS
device design and fabrication, model-based systems engineering, and signal
analysis and processing, including multiphysics modeling, hardware test, and
digital engineering environments. Programmer and data wrangler for
microcontrollers, embedded systems, automated test hardware, workstations,
and cloud-based servers. 

\section{Summary of Skills}

\begin{itemize}

    \item Ph.D. in experimental condensed-matter physics.
    \item Division Chief Scientist, Architecture and System Development Division of the Applied Systems Laboratory (GTRI). 
    \item Model-Based Systems Engineering with OMG SysML (v1 and v2), MagicDraw/Cameo Enterprise Architecture, and FACE conformance tooling.
    \item Open-source developer: author of SysML v2 tooling in Python (sysmlpy, uml2py, sysml\_canvas) with international collaborators.
    \item Innovator and funded author on multiple research, development, and commercialization efforts.
    \item Capable technical writer and presenter.
    \item 25+ years of data acquisition and analysis using Python, C/C++, FORTRAN (77/90/95), Go, Groovy, Java, LabVIEW.
    \item Adept user (and admin) of Unix (Linux, Solaris) and Windows environments, office software suites, and SQL and Graph databases.

\end{itemize}

\section{Experience and Accomplishments}

\begin{topic}

\item [January 2024 - Present] Division Chief Scientist, Architecture and
System Development Division, Applied Systems Laboratory, Georgia Tech
Research Institute, Huntsville, AL.

\begin{itemize}
\item Technical direction, review, and mentoring across division programs in
model-based systems engineering, architecture, and system development.
\end{itemize}

\item [September 2019 - Present] Principal Research Scientist, Architecture
and System Development Division, Applied Systems Laboratory, Georgia Tech
Research Institute, Huntsville, AL.

\begin{itemize}
\item Model Based Systems Engineering technical analyses for US Army Aviation
Future Vertical Lift systems and software architecture.
\item GTRI Technical Lead for Future Long-Range Assault Aircraft (FLRAA)
architectural risk reduction efforts, utilizing the Future Vertical Lift
Architectural Framework (FAF) to enable digital engineering of variant
rotorcraft architectures and incorporating the work of multiple Tier 1, 2,
and 3 corporate architectural partners.
\item Division Chief Scientist of the Architecture and System Development
Division, Applied Systems Laboratory, providing technical direction, review,
and mentoring across division programs.
\item Participating Member of the Architecture Collaboration Working Group
(ACWG) and participant in multiple sub-groups.
\item Jython and Groovy scripting for data extraction and addition into Cameo Enterprise
Architecture.
\item FACE data modeling with the FAME Awesum tool by Tucson Embedded.
\item Open-source development of SysML v2 tooling: clean-room Python
implementations of the UML 2.5.1 metamodel and SysML v1-to-v2 transformation
(uml2py), a SysML v2 reference API and conformance toolchain (sysmlpy, 123/123
OMG XPect conformance), an interactive diagramming canvas (sysml\_canvas), and
a SysML v2 model of the AMS Government Reference Architecture.
\item Clean-room Python ports of the FACE Conformance Test Suite data tooling,
validated byte-for-byte against the reference Java tools.
\end{itemize}

\item [March 2010 - August 2019]System and Software Architect, Project Manager,
Applied Scientist, General Atomics Electromagnetic Systems Group, Huntsville, AL.

\begin{itemize}
\item Modeling of arresting gear system, electromagnetic railgun system, and
laser effect system using the OMG Systems Modeling Language (SysML).
\item Modeling included: use cases, requirements, block definition and internal
block diagrams, requirements diagrams, and state machine diagrams using No
Magic's MagicDraw. Additional MagicDraw functionality added using Jython and
Swing GUI.
\item Integration of SysML architectural models with Windchill product
lifecycle management tool.
\item Lead modeling and simulation of proposed linear displacement transducer
systems for requirements verification and design downselect.
\item Analysis and modeling of pressure data for real-time health prognostication.
\item Development of automated generation of software design documentation from
SysML models using Velocity Templating Language (VTL) and Ruby scripting.
\item Optimization of hypersonic flight trajectories with NASA's OTIS tool.
\item Information and network security analysis with CompTIA Security+
certification.
\item Program and functional management role for remaining SETA work with
AMRDEC/CCDC AMC until Mar. 2018.
\item Provided science and engineering support to the Charles Bowden Weapons
Sciences Laboratory of AMRDEC-WDI, leading in the development of MEMS and
nanomaterials for nano-composite piezoelectric microphones, surface
micro-machined aluminum micro-mirrors, low-pressure microscale sub-die
pressure transducers for hardware health monitoring, surface acoustic health
monitoring of die and printed wiring boards, magnetic micro-actuators for
missile CAS, and MEMS vibrational energy health transducers for missile
health prognostication.
\item Functional manager of employees supporting work at the Weapons Sciences
Laboratory.
\end{itemize}

\item [July 2007 - March 2010]Thin Film Process Engineer, Gulton, Inc.,
South Plainfield, NJ.

\begin{itemize}
    \item Successfully developed thin film thermal print head (TPH) product lines for a small company with no corporate thin film manufacturing experience.
    \item Process control and product management of existing thick film TPHs from screening to final assembly.
    \item Selection, installation, maintenance of production physical vapor deposition equipment.
    \item Developed dry and wet etch processes for metals, resistor materials, and passivation dielectrics from prototype to production.
    \item Programmed automated test equipment for resistor testing and trimming at the die level and the packaged part level.
    \item Developed and implemented web application servers for SPC data collection and analysis across the factory floor.
\end{itemize}

\item [April~2000-July 2007]Principal Scientist, Research Support Instruments,
Princeton, NJ.

\begin{itemize}
    \item Conceived and managed numerous technology development projects in micro-fabricated structures.
    \item Appointed Visiting Professional Technical Staff Member II in the Department of Mechanical and Aerospace Engineering of Princeton University (2000--2006).
    \item Industrial Partner with Princeton Institute of the Science and Technology of Materials.
    \item Awarded multiple research grants towards RSI micro-machining efforts from the SBIR/STTR program with AFRL, ONR, NIST, \& NASA.
    \item Principal Investigator/Project Manager, NASA Small Business Innovation Research (SBIR) Grant Phase I ``An inexpensive, versatile micro-machined mass spectrometer.''
    \item Principal Investigator/Project Manager, Air Force SBIR Grant Phase I \& II, ``A micro-machined energy-mass spectrograph for micro-satellite applications.''
    \item Principal Investigator/Project Manager, NIST SBIR Grant Phase I, ``A Micro-Electroformed Cross Capacitor.''
    \item Principal Investigator/Project Manager, ONR STTR Grant Phase I, ``Passive MEMS tagging of Personnel.''
    \item Modeled the performance of micro-fabricated ion optics using Numeric Python and OpenDX.
\end{itemize}

\item [1993-2000]Research Assistant, Penn State Physics Department, University
Park, PA.

\begin{itemize}

\item Developed skills in many analytical techniques for materials and surfaces:
including Raman Spectroscopy, Photoluminescence, Electron Energy Loss
Spectroscopy (HREELS), Photoelectron Emission Spectroscopies (UPS,
XPS, AES), LEED and RHEED, and powder neutron diffraction.

\item Studied a wide range of novel materials: fullerenes; Si, Ge, and C
amorphous and crystalline nanoparticles; ferromagnetic oxides; and
II-VI epitaxial films.

\item Developed a wide range of high and ultra-high vacuum materials deposition
skills including rf \& dc magnetron sputtering, e-beam evaporation,
and pulsed laser deposition.

\item Authored peer reviewed scientific papers and funding reports for the
National Science Foundation and Department of Energy.

\end{itemize}
\end{topic}

\section{Education }

\begin{topic}
\item [Ph.~D.~Physics]The Pennsylvania State University, University Park,
PA, May 2000. Thesis: ``\emph{in situ} studies of the electronic and vibrational
properties of thin films and novel materials'', Advisor: the late Professor
Jeff Lannin. Thesis completed under the supervision of committee chair
Professor Robert Collins.
\item [B.S.~Physics]\emph{Magna~cum~laude.} University of Missouri-Rolla,
Rolla, MO, 1993. Minor in mathematics.
\end{topic}

\section{Current Memberships }

\begin{itemize}
\item IEEE/ Computer Society
\item the International Council on Systems Engineering (INCOSE) 
\end{itemize}

\section{Current Certifications}

\begin{itemize}
\item INCOSE Certified Systems Engineering Professional (CSEP), March 2025
\end{itemize}

\section{Open Source \& Publications}

\begin{itemize}
\item Author and maintainer of the sysmlpy Python toolchain for the OMG SysML
v2 standard (parser, semantic analyzer, API services), with international
collaborators and OMG conformance testing. \texttt{github.com/mycr0ft}
\item Curator of the \emph{awesome-sysml} list of SysML v2 language tools.
\end{itemize}

\begin{thebibliography}{1}

\bibitem{horowitz2013} Stephen B. Horowitz, Michael S. Allen, Jon R. Fox, Jean
P. Cortes, Laura Barkett, Adam D. Mathias, Corey Hernandez, Adam C. Martin,
Mohan Sanghadasa, Stephen Marotta ``A low frequency MEMS vibration sensor for
low power missile health monitoring.''  2013 IEEE Aerospace Conference, Big
Sky, MT,  March 2013.

\bibitem{fox2012mirrors} Jon R. Fox, Adam D. Mathias, Jean Cortes, Michael S.
Allen, Stephen B. Horowitz,  Mark G. Temmen, and Mohan Sanghadasa.
``Fabrication and characterization of ultra-fast electrostatically-actuated
surface micro-machined aluminum mirrors.'' American Vacuum Society 59th Annual
International Symposium and Exhibition held in Tampa, Florida on October 30,
2012.

\bibitem{fox2012microphones}  J.R. Fox, S.B. Horowitz, J.P. Cortes, M.S. Allen,
A.D. Mathias, L.A. Barkett, and M. Sanghadasa. ``Fabrication and Testing of
Suspended Piezoelectric Nanocomposite Membranes.'' American Vacuum Society 59th
Annual International Symposium and Exhibition held in Tampa, Florida on October
29, 2012.

\bibitem{horowitz2012}
Stephen Horowitz, Jon Fox, Dustin Mathias, Jean Cortes, Michael Allen, Laura Barkett, Mohan Sanghadasa. ``Miniaturized Nanocomposite Piezoelectric Microphones for UAS Applications.'' 2012 Joint Meeting of the Military Sensing Symposia (MSS) held in Washington, DC on October 22-25, 2012.

\bibitem{horowitzSA2012} Stephen Horowitz, Dustin Mathias, Jon Fox, JP Cortes, Mohan Sanghadasa, Paul Ashley.  ``Effects of Scaling and Geometry on the Performance of Piezoelectric Microphones.'' Sensors and Actuators, 185: 24-32, October 2012.

\bibitem{mathias_2011} Adam D Mathias, Jon R Fox, Jean P Cortes, Stephen B Horowitz, Mohan Sanghadasa, Paul Ashley. ``Recent progress toward a nanostructured piezoelectric microphone.'' The Journal of the Acoustical Society of America Volume 130 (4) 2427, April 2011.

\bibitem{mathias_2011infotech} Adam D Mathias, Jon R Fox, Stephen B Horowitz, Mark G Temmen. A High Speed Electrostatically Actuated Al-MEMS Micromirror Array. Infotech@ Aerospace 2011, March 29, 2011.

\bibitem{salamon} M. M. Salamon, J. R. Fox, D. J. Sullivan, J. F. Kline,  ``Fabrication of a novel suspended microactuator MEMs device.'', Society of Manufacturing Engineers MicroManufacturing Conference,  Minneapolis, Minnesota, Mar. 2009.

\bibitem{mcandrew} McAndrew, B., Kline, J., Fox, J.R., Sullivan, D.J., Miles,
R., \textit{``Supersonic Vehicle Control by Microwave Driven Plasma
Discharges,''} AIAA--2002--0354, 40th Aerospace Sciences Meeting \& Exhibit,
Reno, Nevada, Jan. 14--17, 2002.

\bibitem{kline_2001} Kline, J., J.R. Fox, V. DiPietro, and J. Brandenburg.  A Review
of the Compact Micro-Torus eXperiment. \emph{4th Symposium on the
Current Trends in International Fusion Research: A Review.} Washington,
D.C. 2001.

\bibitem{kline_2000} John F. Kline, Jon R.
Fox, Vincent DiPietro, John Vandenburg, Characterization of the Compact
Micro Torus eXperiment. The 27th IEEE International Conference on Plasma
Science. New Orleans, LA. June 6, 2000.

\bibitem{kline_nasa_workshop} Kline, J., J.R. Fox, V. DiPietro, and J.
Brandenburg.  Propulsion Applications of the Compact Micro-Torus
eXperiment. \emph{NASA Fusion Propulsion Workshop.} 2000.

\bibitem{sirenko_2000} A. Sirenko, J. R. Fox, I. A. Akimov, X.X. Xi, S. Rouvimov, and Z.
Lilienthal-Weber. \emph{In situ} Raman scattering studies of amorphous
and crystalline states in Si nanoparticles. \emph{Solid State Commun.},
113:553, 2000.

\bibitem{xi_2000} X. X. Xi, H. C. Li, W. D. Si, A. A. Sirenko, I. A. Akimov,
J. R. Fox, A. M. Clark, J. H. Hao. Oxide thin films for tunable microwave
devices. \emph{Journal of Electroceramics}, 4 (2-3): 393-405, 2000.

\bibitem{fox_1998} J. R. Fox, I.A. Akimov, X.X. Xi, and A. A. Sirenko. In situ studies
of electronic and vibrational properties of Si nanoparticles. In L.T.
Canham, \emph{et al.} , editors, \emph{Advances in Microcrystalline
and Nanocrystalline Semiconductors}, volume 536, page 287. Materials
Research Society, 1998.

\bibitem{fox_1996}J. R. Fox, G. P. Lopinski, and J. S. Lannin, G. B. Adams, J. B. Page,
and J. E. Fischer. Neutron diffraction and structural models of RbC\( _{60} \)
phases. \emph{Chem. Phys. Lett.}, 249(3):195, 1996.

\bibitem{lopinski_1996} G.P. Lopinski, J.R. Fox, J.S. Lannin, F.S. Flack, and N. Samarth.
Reconstruction induced changes in the electronic states of ZnSe(100).
\emph{Surf. Sci. Lett.}, 355(1/3):L355, 1996.

\bibitem{lopinski_1995} G. P. Lopinski, J. R. Fox, and J. S. Lannin. Electronic and vibrational
properties of laser modified C\( _{60} \). \emph{Chem. Phys. Lett.}
239(1/3):107-111, 1995.

\bibitem{lopinski_1994} G. P. Lopinski, M.G. Mitch, J. R. Fox, and J. S. Lannin. Phase transition
in RbC\( _{60} \): Vibrational and electronic properties. \emph{Phys.
Rev. B.}, Rapid Commun., 50(21): 16098-16101, 1994.

\end{thebibliography}

\end{document}
